\begin{recipe}{Quiche met groente in een cirkel}
\kicker[Hoofdgerecht,Vegetarisch]{Hoofdgerecht · vegetarisch}
\index[register]{Quiche met groente in een cirkel}
\index[register]{Courgette}
\index[register]{Aubergine}
\index[register]{Geitenkaas}

\meta{100 minuten \dvd Voor 5 porties \dvd Oven 180\,°C}

\lettrine{E}{en} quiche waarbij de groente van buiten naar binnen in cirkels worden
gelegd, wat een mooi visueel effect geeft aan tafel. De valkuil is dat de groente
rauw blijft.

\blockrule{Ingrediënten}

\noindent\textit{Deeg}
\begin{ingredients}
  \ing{150 g}{bloem}
  \ing{75 g}{ongezouten roomboter, zacht}
  \ing{1}{ei}
  \ingb{zout}
\end{ingredients}

\medskip
\noindent\textit{Vulling}
\begin{ingredients}
  \ing{200 g}{geitenkaas schijfjes naturel}\index[register]{Geitenkaas}
  \ing{3}{winterpenen}\index[register]{Wortel}
  \ing{1}{courgette}\index[register]{Courgette}
  \ing{1}{aubergine}\index[register]{Aubergine}
  \ing{1 flinke lepel}{Griekse yoghurt}
  \ing{2}{eieren}
  \ingb{verse tijm}
  \ingb{knoflookpoeder, uienpoeder, versgemalen peper, zout}
\end{ingredients}

\blockrule{Bereiding}
\tip{Snijd de groente zo dun mogelijk, bij voorkeur met een kaasschaaf. De oven
gaart ze nauwelijks verder, dus blancheer goed voor.}
\begin{steps}
  \item Maak het deeg: kneed de zachte boter met de bloem, het ei en een snuf zout
        tot een samenhangend geheel. Wikkel in vershoudfolie en laat minstens 15 minuten
        rusten in de koelkast.
  \item Schil de wortels. Was de courgette en aubergine. Snijd alle groente in dunne
        plakken, bij voorkeur met een kaasschaaf.
  \item Breng een pan gezouten water aan de kook en blancheer alle groente 2 minuten.
        Giet af, laat uitlekken en dep de plakken goed droog met keukenpapier:
        tussen dunne plakken blijft veel water hangen, en dat komt er in de
        oven weer uit.
  \item Meng de geitenkaas, Griekse yoghurt en eieren tot een egaal mengsel. Breng
        flink op smaak met knoflookpoeder, uienpoeder, peper en zout. Niet stijf
        kloppen.
  \item Verwarm de oven voor op 180\,°C (boven- en onderwarmte). Vet de springvorm in, strooi er een beetje
        bloem in en klop de vorm rond zodat de bloem de bodem en wanden bedekt.
  \item Rol het deeg uit en bekleed de springvorm. Prik met een vork wat
        gaatjes in de bodem: dat voorkomt dat het deeg tijdens het bakken
        opbolt. Leg de groenteplakken van de
        buitenrand naar het midden in concentrische cirkels, afwisselend per soort.
  \item Giet de vulling erover en schud de vorm licht zodat de vulling ook tussen de
        groente terechtkomt.
  \item Bak 45 tot 50 minuten op 180\,°C: de custard is hier dunner dan bij
        de andere quiches, dus vul de vorm niet te hoog en gun hem de extra
        tijd als hij nog golft. Laat de quiche 10 minuten rusten voor het
        aansnijden en strooi dan vlak voor het serveren wat losse tijmblaadjes
        erover.
\end{steps}

\heroimagefade{images/quiche-groente-cirkel}
\end{recipe}
